\chapter{Radon measures}\label{sec:radon-measures}
Weak convergence is useful because geometric limits need not remain smooth. We first specify the measures and the topology used in the later chapters. The reference is \cite{Sim18}, Chapter 1, Sections 1 and 5.

\begin{definition}[Radon measure]\label{def:radon-measure}
A nonnegative Borel measure $\mu$ on an open set $U\subset\mathbb{R}^n$ is a \emph{Radon measure} if it is finite on compact subsets of $U$ and, for every Borel set $A\subset U$,
\[
 \mu(A)=\inf\{\mu(O):A\subset O\subset U,\ O\text{ open}\}
       =\sup\{\mu(K):K\subset A,\ K\text{ compact}\}.
\]
We use its completion when integrating measurable functions. Its support consists of the points $x\in U$ for which every neighborhood has positive measure. Local finiteness does not require $\mu(U)<\infty$. See \cite{Sim18}, Chapter 1, Section 5, Definition 5.6, p.~31, for the corresponding outer-measure formulation.
\end{definition}

\begin{definition}[Weak convergence of measures]\label{def:weak-measure}
Radon measures $\mu_j$ converge \emph{weakly in $U$} to a Radon measure $\mu$, written $\mu_j\rightharpoonup\mu$, if
\[
 \lim_{j\to\infty}\int_U\varphi\,d\mu_j=\int_U\varphi\,d\mu
 \qquad\text{for every }\varphi\in C_c(U).
\]
This is also called vague convergence. Testing only compactly supported functions allows mass to escape toward infinity or the boundary of $U$. See \cite{Sim18}, Chapter 1, Section 5.
\end{definition}

\begin{theorem}[Local compactness of Radon measures]\label{thm:radon-compactness}
Let $\mu_j$ be nonnegative Radon measures on $U$ and suppose that
\[
 \sup_j\mu_j(K)<\infty\qquad\text{for every compact }K\subset U.
\]
There are a subsequence $\mu_{j_k}$ and a Radon measure $\mu$ on $U$ such that $\mu_{j_k}\rightharpoonup\mu$ in $U$.

See \cite{Sim18}, Chapter 1, Section 5, Theorem 5.15, pp.~38--39. The Euclidean countable-test argument and positive representation theorem are developed below.
\end{theorem}
\begin{proof}
Lemma \ref{lem:measure-diagonal-functional} extracts a subsequence and a positive linear functional $L$ such that $\int f\,d\mu_{j_k}\to L(f)$ for every $f\in C_c(U)$. Theorem \ref{thm:positive-riesz-representation} gives a Radon measure $\mu$ with $L(f)=\int f\,d\mu$. This is exactly the convergence in Definition \ref{def:weak-measure}. The countable family exists because $U$ is a Euclidean open set; we do not infer sequential extraction from weak-star compactness on a general nonmetrizable space.
\end{proof}

\begin{proposition}[Lower semicontinuity on open sets]\label{prop:measure-lsc}
If $\mu_j\rightharpoonup\mu$ in $U$, then
\[
 \mu(O)\leq\liminf_{j\to\infty}\mu_j(O)
 \qquad\text{for every open }O\subset U.
\]
This is the open-set inequality accompanying weak convergence in \cite{Sim18}, Chapter 1, Section 5.
\end{proposition}
\begin{proof}
Fix $0\leq\varphi\leq1$ with $\varphi\in C_c(O)$. Weak convergence gives
\[
 \int\varphi\,d\mu=\lim_j\int\varphi\,d\mu_j\leq\liminf_j\mu_j(O).
\]
For every compact $K\subset O$ we can choose such a function with $\varphi=1$ on $K$. Taking the supremum over $K$ and using inner regularity proves the assertion, including the case $\mu(O)=\infty$.
\end{proof}

\chapter{Hausdorff measure and density}\label{sec:hausdorff-measure}
Hausdorff measure gives a notion of $m$-dimensional size before any tangent plane is known. We use the normalization in \cite{Sim18}, Chapter 1, Section 2.

\begin{definition}[Normalized Hausdorff measure]\label{def:hausdorff-measure}
Let $m\geq1$ be an integer and $A\subset\mathbb{R}^n$. For $\delta>0$ define
\[
 \mathcal{H}^m_\delta(A)=\inf\left\{
 \sum_{i=1}^{\infty}\omega_m\left(\frac{\operatorname{diam}E_i}{2}\right)^m:
 A\subset\bigcup_iE_i,\quad\operatorname{diam}E_i<\delta\right\},
 \qquad
 \mathcal{H}^m(A)=\lim_{\delta\downarrow0}\mathcal{H}^m_\delta(A).
\]
The empty cover gives value zero for $A=\varnothing$. The limit exists in $[0,\infty]$ by monotonicity as $\delta$ decreases. We set $\mathcal{H}^0$ equal to counting measure. These are Borel regular outer measures. With this normalization $\mathcal{H}^m=\mathcal{L}^m$ on $\mathbb{R}^m$; that identification is a cited result, not part of the definition. See \cite{Sim18}, Chapter 1, equations (2.1)--(2.5) and Theorem 2.6, pp.~11--14, and \cite{Fed69}, Section 2.10.2(1), p.~171.
\end{definition}

\begin{proposition}[Lipschitz distortion of Hausdorff measure]\label{prop:lipschitz-measure}
Let $A\subset\mathbb{R}^n$, let $f:A\to\mathbb{R}^q$ be $L$-Lipschitz with $0<L<\infty$, and let $m\geq1$. Then
\[
 \mathcal{H}^m(f(A))\leq L^m\mathcal{H}^m(A).
\]
In particular, Lipschitz maps take $\mathcal{H}^m$-null sets to $\mathcal{H}^m$-null sets. A constant map also has this null-set property. This follows directly from the covering definition in \cite{Sim18}, Chapter 1, Section 2.
\end{proposition}
\begin{proof}
Intersect each set in a $\delta$-cover of $A$ with $A$ and take its image under $f$. The image diameters are at most $L$ times the original diameters. Thus
\[
 \mathcal{H}^m_{L\delta}(f(A))\leq L^m\mathcal{H}^m_\delta(A).
\]
Let $\delta\downarrow0$. For a constant map the image is empty or a singleton, whose $\mathcal{H}^m$ measure is zero when $m\geq1$.
\end{proof}

\begin{definition}[Density]\label{def:density}
For a Radon measure $\mu$ on $U$, a point $x\in U$, and $m\geq1$, define
\[
 \Theta^{*m}(\mu,x)=\limsup_{r\downarrow0}\frac{\mu(B_r(x))}{\omega_m r^m},
 \qquad
 \Theta_*^m(\mu,x)=\liminf_{r\downarrow0}\frac{\mu(B_r(x))}{\omega_m r^m}.
\]
When these agree, their common value is the density $\Theta^m(\mu,x)$. The values may be infinite. For $\mu=\mathcal{H}^m\llcorner M$ we also write $\Theta^m(M,x)$. See \cite{Sim18}, Chapter 1, Section 3, p.~14.
\end{definition}

\begin{example}[A plane with multiplicity]\label{ex:plane-density}
Let $P$ be an affine $m$-plane, let $\theta>0$, and set $\mu=\theta\mathcal{H}^m\llcorner P$. For every $x\in P$ and $r>0$,
\[
 \mu(B_r(x))=\theta\omega_mr^m,\qquad \Theta^m(\mu,x)=\theta.
\]
Indeed, a Euclidean isometry identifies $P\cap B_r(x)$ with an $m$-dimensional ball, and the normalized Hausdorff measure agrees with Lebesgue measure on that plane. This model fixes the normalization used for tangent planes and varifolds below.
\end{example}
